\subsection{Summary}
\label{comesh:section:conclusion}

\sysname{} demonstrates that a fully-decentralized Sense-Trigger-Actuation control plane for large commercial edge IoT meshes is both theoretically sound and practically deployable.
Raspberry Pi deployments show that \sysname{} reduces per-device CPU load by over $6\times$ compared to a three-hub baseline, cuts inter-device bandwidth by $10\times$ versus a central hub, and achieves median coterie operations (quorum reply, leader election, state transfer) in under 2.1 seconds.
Large-scale simulations confirm that \sysname{} maintains stable latencies under up to 20\% device churn and 5\% packet loss, and scales to hundreds of devices without degradation.
Today's emerging IoT standards---Matter~\cite{csa2025matter}, Thread~\cite{thread2020fundamentals}, and others---align with \sysname's vision of a truly autonomous, cloudless control plane for mesh-based edge environments.
Future work can explore Byzantine faults, extreme mobility (e.g., vehicular networks), and city-scale deployments.
\rev{Highly mobile deployments---e.g., large factories where humanoid robots move between stations and roam the floor---are a particularly interesting target: \sysname{} can already absorb such movement \emph{reactively} through its churn- and mobility-handling coterie migration, but predicting robot motion (e.g., from task schedules or learned trajectories) would let \sysname{} \emph{proactively} reshape coteries before a robot moves, keeping coordination local to wherever the robot is about to be.}
